\chapter{Classical Mechanics}
\label{cm}

{
\setlength\epigraphwidth{0.6\textwidth}
\epigraph {
    Motion\\
There is no motion, said a wise old Greek.\\
The other sage just smiled and walked around.\\
A better answer hardly could be found,\\
And everybody praised this tongue in cheek.\\
But, gentlemen, this funny story might\\
Remind us of a man of great renown:\\
Day after day the Sun walks up and down,\\
And yet the stubborn Galilei was right.
} {A.S. Pushkin, Translated by D. Broytman }
}

\section{The Two Roles of Classical Mechanics}
\label{cm:tr}

One of the most important models in physics is mechanics.
Until the end of the 19th century, mechanics was considered as direct and the only
the basis of all physics. By words understand or
to explain a physical phenomenon meant constructing a mechanical model of it.

Situation has changed,  other physical models exist, e.g. quantum mechanics so 
"usual" mechanics has changed its name to "classical mechanics" and still is
the central part of theoretical physics. All main physical concepts emerged from
it.


Classical mechanics occupies a peculiar position among physical theories. It is,
first, a complete and self-contained subject: a body of concepts, equations, and
results describing the motion of bodies under forces, valid and useful on its own
terms, with no reference whatsoever to the Planck constant $\hbar$ or the speed of
light $c$. Celestial mechanics, rigid-body dynamics, the theory of small
oscillations, continuum mechanics, and the modern theory of dynamical systems and
chaos are all classical mechanics pursued for its own sake, and none of them is
waiting to be corrected by a more fundamental theory in order to be meaningful.

Second, and for the purposes of this book more importantly, classical mechanics is
the \emph{base} on which relativistic and quantum mechanics are built. Special and
general relativity are commonly introduced as corrections to Newtonian mechanics
that become significant when velocities approach $c$ or fields become strong;
quantum mechanics is commonly introduced as a modification that becomes significant
at the scale of $\hbar$. In both cases, the classical theory is treated as a known,
well-understood limit --- the solid ground from which the generalization departs
and to which it must reduce.

It is tempting to think that once this second, subordinate role is granted, nothing
further needs to be said about classical mechanics: its job is simply to sit still
and be a limit. The purpose of this chapter, and of the chapters immediately
following it, is to argue that this is false in two related ways.

\begin{enumerate}
  \item Even restricted to the narrow task of serving as the base for
  relativistic and quantum generalization, classical mechanics is not a settled,
  transparent starting point. Producing the structure that gets generalized ---
  a well-defined phase space, a Hamiltonian, a consistent variational principle
  --- is itself a nontrivial achievement that fails, or becomes ambiguous, for
  entire classes of ordinary classical systems.
  \item Several phenomena that are routinely presented as characteristically
  relativistic or characteristically quantum-mechanical --- an underdetermination
  of physical ontology by the observable facts, an ambiguity in the passage from
  classical description to quantized theory, a geometric ``phase'' with no
  classical counterpart, an apparent conflict between microscopic reversibility
  and macroscopic causal, dissipative behavior --- already occur inside classical
  mechanics, before relativity or $\hbar$ ever enter the discussion. When such a
  paradox survives the excision of everything relativistic or quantum, its source
  cannot honestly be attributed to relativity or to quantum theory.
\end{enumerate}

\section{Classical Mechanics as a Subject in Its Own Right}
\label{cm:st}

It is worth pausing on the first role only long enough to insist that it is real.
The three-body problem, the precession of orbits, the stability of the solar
system, the dynamics of rigid bodies and gyroscopes, the propagation of waves in
elastic media, the onset of turbulence, and the sensitive dependence on initial
conditions studied by dynamical-systems theory are all substantial, still
partially open, areas of physics conducted entirely within the classical
framework. None of this material is a mere warm-up exercise for quantum mechanics,
and a physicist could spend a full career inside it without ever needing $\hbar$.
This book is not about that role. It is mentioned here only to make clear that
when the rest of this book treats classical mechanics as a base to be generalized,
it is not thereby claiming that this exhausts what classical mechanics is.

The four cases below are developed fully in later chapters of this book. They
are introduced here, briefly and together, because their common moral is best
seen by contrast: in each case, a feature usually advertised as a signature of
quantum theory or of relativity turns out to have an exact, or near-exact,
analogue that requires nothing beyond Newton, Lagrange, and Hamilton.

\subsection{Hertz's forceless mechanics }

In his \emph{Principles of Mechanics}, Heinrich Hertz reformulated dynamics
without the concept of force at all. Observable bodies are treated as coupled to
unobservable ``hidden masses'' through rigid, holonomic constraints, and the
combined system moves along a geodesic of a metric built from the mass
distribution on an enlarged configuration space. Suitably arranged, Hertz's
constrained geodesic motion reproduces the trajectories that Newtonian mechanics
attributes to the action of a force. The two formulations are empirically
equivalent: no experiment performed on the observable bodies alone can decide
between ``a force is acting'' and ``a constraint is enforcing geodesic motion in
a larger space.''

This is a case of underdetermination of theory --- indeed of ontology --- by
observation, arising entirely within classical mechanics. It is normally quantum
mechanics that is credited with forcing physicists to confront multiple,
empirically indistinguishable accounts of what is ``really'' happening beneath
the observed statistics: Bohmian trajectories, Everettian branches, and
epistemic and objective-collapse readings of the wavefunction are all consistent
with the same experimental record. Hertz shows that the same structural
situation --- one set of observations, several inequivalent stories about the
underlying ontology --- already existed in nineteenth-century particle mechanics,
with no probability amplitude in sight.

\subsection{Nonholonomic constraints }

A constraint of the Pfaffian form
\begin{equation}
  \sum_i a_i(q)\,\dot q_i = 0
\end{equation}
is \emph{nonholonomic} when it is not integrable to a relation $f(q)=0$ among the
coordinates alone --- the standard example being a disk or coin rolling without
slipping on a plane. Such systems present a genuine fork in the road. Applying
d'Alembert's principle directly, with the constraint enforced at the level of
virtual displacements, yields the correct equations of motion, confirmed by
experiment. Applying Hamilton's variational principle instead --- extremizing the
action with the constraint imposed by a Lagrange multiplier, as one legitimately
may for holonomic constraints --- yields a different, and physically incorrect,
set of equations. This is the long-standing nonholonomic-versus-vakonomic
controversy: two superficially reasonable derivations from the same starting
Lagrangian and the same constraint disagree about the dynamics, and only one of
them is right.

One consequence is that nonholonomic systems generically fail to admit a
straightforward canonical Hamiltonian formulation on the naive phase space:
the constrained bracket structure one is forced to use is not the canonical
Poisson bracket, and which reduced structure is the correct one to use is itself
a nontrivial derivation, system by system. This matters directly for the
generalization picture of Section~\ref{cm:ba}: canonical quantization
needs a Poisson bracket to promote to a commutator, and here the classical
bracket is already ambiguous before any promotion is attempted. The
operator-ordering problem --- the fact that a classical observable such as $qp$
does not determine a unique quantum operator, since $\hat q\hat p$, $\hat p\hat
q$, and $\tfrac12(\hat q\hat p+\hat p\hat q)$ all reduce to it as $\hbar\to0$
--- is usually presented as an intrinsically quantum indeterminacy. In the
nonholonomic case the indeterminacy has already begun classically, in the
disagreement over what the correct pre-quantization bracket even is.

\subsection{Shape-changing bodies and classical geometric phase}

A body that changes its own shape --- the paradigm case is a cat, initially
oriented feet-up, that reorients itself in mid-air with zero net angular
momentum by cyclically changing its shape --- cannot be treated by rigid-body
mechanics. One instead works on a \emph{shape space} $\mathcal S$ (the space of
internal configurations, with rigid motions quotiented out) and finds that the
conservation of angular momentum, together with the equations of motion,
defines a connection on a principal bundle over $\mathcal S$ with structure
group $SO(3)$. A closed loop in shape space --- the cat returns to its initial
shape --- need not be a closed loop in the full configuration space: the body
can acquire a net rotation purely as a holonomy of this connection, with no
external torque and no change in angular momentum at any instant. The same
mechanism governs low-Reynolds-number swimmers and, more generally, any system
in which a periodic change of internal, ``slow'' variables produces a net shift
in an external, cyclic variable; the resulting angle is called the Hannay angle.

The geometric or Berry phase acquired by a quantum system whose Hamiltonian is
slowly and cyclically varied is frequently introduced as a paradigmatically
quantum-mechanical effect, precisely because it has no counterpart in the naive
classical picture of a particle following a trajectory. The falling cat and the
Hannay angle show that this is too strong a claim about geometric phases in
general: the underlying mechanism --- holonomy of a connection around a loop in
parameter or shape space, decoupled from the local dynamics along the way --- is
already fully present, and was fully worked out, in ordinary classical
mechanics. What is genuinely new in the quantum case is the specific bundle and
connection defined by the adiabatic theorem applied to a Hilbert space of
states; the geometric phenomenon of holonomy itself is not new.

\subsection{Open systems}

Two classical difficulties are worth setting side by side. The first is
statistical: Boltzmann's $H$-theorem derives macroscopic irreversibility from
microscopically time-reversible mechanics, and was met immediately by
Loschmidt's reversibility objection (reversing all velocities should reverse
the approach to equilibrium) and Zermelo's recurrence objection (Poincar\'e
recurrence guarantees a return arbitrarily close to any earlier state). Both
objections are correct as stated, and reconciling them with the observed arrow
of time is a genuinely hard, and still actively discussed, problem --- entirely
classical, with no reference to measurement or wavefunction collapse.

The second is dynamical: an accelerating point charge radiates, and Newton's
second law must be supplemented by the Abraham--Lorentz self-force,
\begin{equation}
  m\dot v = F_{\text{ext}} + \frac{2}{3}\frac{e^2}{c^3}\,\dot a \, ,
\end{equation}
a third-order correction proportional to the derivative of the acceleration.
This equation admits runaway solutions, in which the particle self-accelerates
exponentially with no external force at all; demanding that solutions remain
bounded, instead, forces the particle to begin accelerating \emph{before} the
external force is applied --- a pre-acceleration that appears to violate
causality, entirely within classical electrodynamics.

Both irreversibility and an apparent tension with causal, deterministic
evolution are frequently treated as problems that quantum theory introduces
--- through the measurement postulate, decoherence, or the collapse of the
wavefunction --- or that relativity introduces, through its reworking of
simultaneity. Loschmidt, Zermelo, and Abraham--Lorentz show that a purely
classical system, governed by Newton's laws and Maxwell's equations and nothing
else, already exhibits dissipative, apparently irreversible, and apparently
acausal behavior that has never been fully tamed.

\subsection{Looking Ahead}

Each of the four cases above is the subject of a full chapter later in this
book (chapters \ref{hm},\ref{nh},\ref{shape} and \ref{open})
, where the mathematics is developed in detail and the historical debate is
traced more carefully. The point of gathering them here, in outline, is
methodological: before a ``paradox'' encountered in relativity or in quantum
mechanics is attributed to relativity or to quantum mechanics as such, it is
worth checking whether classical mechanics, unaided, already contains a version
of it. Where it does, what is new in the generalization is not the paradox but,
at most, its dressing --- and the genuinely new content of the generalization
has to be located elsewhere.

\section{Classical Mechanics as the Base of Generalization}
\label{cm:ba}

The textbook picture of generalization is familiar. Quantization is presented as
a recipe applied to a classical Hamiltonian system: promote canonical coordinates
$(q,p)$ to operators $(\hat q,\hat p)$, and promote the Poisson bracket to the
commutator,
\begin{equation}
  \{f,g\} \;\longmapsto\; \frac{1}{i\hbar}\,[\hat f,\hat g] \, .
\end{equation}
The inverse limit, $\hbar \to 0$, recovers classical mechanics via the WKB
approximation and Ehrenfest's theorem. Similarly, special relativity is presented
as reducing to Newtonian mechanics for $v \ll c$, and general relativity as
reducing to Newtonian gravitation in the weak-field, slow-motion limit.

Both halves of this picture presuppose that classical mechanics has already
delivered, for the system in question, a specific and unambiguous piece of
structure: a phase space with a well-defined symplectic (or Poisson) structure,
a Hamiltonian generating the correct equations of motion, and a family of
observables closed under the bracket. Only given that structure does
``quantize it'' or ``take the classical limit'' pick out a definite theory.

For a mass point moving under a potential, or for a system subject only to
holonomic constraints, this presupposition is satisfied and the textbook picture
is accurate. But it is easy to overstate how general this good case is. As soon
as one leaves it --- constraints that cannot be integrated to relations among
coordinates alone, systems that exchange energy or momentum with an unmodeled
environment, systems whose configuration space itself is not fixed once and for
all --- the classical structure needed as input to generalization becomes
ambiguous, multiple, or simply unavailable in closed form. In those cases the
question ``what is the classical mechanics of this system'', asked with no
reference to $\hbar$ or $c$ at all, already has more than one defensible answer.
That is the sense in which classical mechanics, even in its purely subordinate
role as a base, is \emph{far from trivial}: the base itself has to be constructed,
and the construction is not always unique.

As a physical model classical mechanics is based on three principles.
\begin{description}
  \item[ Global Simultaneity] 
    Time is universal and global in classical mechanics. State of the system is determined by
    positions and velocities at fixed time. Every particle has its own coordinates however time
    is common. We can always measure time  exactly and instantly without any influence 
    on mechanical systems. Time is continuous, discrete time would lead to finite difference
    equations instead of differential ones. So dynamics would be undefined, since one always can
    add periodical function to solution.
  \item[ Instant Interaction]
    All interactions are instant i.e. depend on particle states at one and the same moment of time.
    In particular it means that all accelerations and hence evolution is determined by one state (not necessary initial one). It is possible to think that any interaction takes small but finite period of time. This supposition will lead to functional equations instead of differential ones.
  \item [ Precision]
    We can find out and set up exact value of any variable at any time instantly without disturbing
    the system in any way. In particular we can set exactly positions and velocities which means 
    among other things that space is continuous.
\end{description}

There is a great number of books about classical mechanics 
(see e.g\cite{Arnold-cm},\cite{Landau-Lifshitz-1},
\cite{Lanczos},\cite{Goldstein}). 

In Ideal world of classical mechanics space is three-dimensional  and euclidean
$\mathbb{R}^3$, mass is real number and time is one-dimensional $\mathbb{T}$.
So universe of classical mechanics is $\mathbb{T} \times \mathbb{R}^3$. 

 The  group of galilean transformations acts on universe \cite{Arnold-cm}. 
Every galilean transformation can be written in a unique way as the composition of a rotation, a translation,
and a uniform motion ($g = g_1 \circ g_2 \circ g_3$), where
\[g_1(t,\vec x)=(t,\vec x+\vec v t)\qquad \forall t\in \mathbb{R},\vec x \in\mathbb{R}^3\] is uniform motion with velocity $\vec v \in\mathbb{R}^3$,
\[g_2(t,\vec x)=(t+t_0,\vec x+\vec x_0)\qquad \forall t\in \mathbb{R},\vec x \in\mathbb{R}^3\] is translation of the origin, and
\[g_3(t,\vec x)=(t,O\vec x)\qquad \forall t\in \mathbb{R},\vec x \in\mathbb{R}^3\] rotation of the coordinate axes ($O:\mathbb{R}^3 \rightarrow \mathbb{R}^3$ is an orthogonal transformation).

Let $M$ be a set. A one-to-one correspondence 
\[\phi_1: M \rightarrow \mathbb{R} \times \mathbb{R}^3\]
is called a galilean (inertial) coordinate system on the set $M$. A coordinate system $\phi_2$ moves
uniformly with respect to $\phi_1$  if $\phi_1\circ\phi_2^{-1}$ is a galilean
transformation. 

Inertial coordinate systems  possess the following two properties:
\begin{enumerate}
    \item All the laws of nature at all moments of time are the same in all inertial coordinate systems;
    \item  All coordinate systems in uniform rectilinear motion with respect to an inertial one are themselves inertial.
\end{enumerate}
This is called `Galileo's principle of relativity'.

The simplest object of classical mechanics is point mass (particle) i.e. object whose shape, 
sizes and internal structure are inessential, so mass (positive scalar $m$) is the only parameter. 
The state of a particle is fully determined by position ($\vec x \in\mathbb{R}^3$) 
and velocity ($\vec v \in\mathbb{R}^3$). 
More accurately: A motion in $\mathbb{R}^n$ is a differentiable mapping 
$\vec x: \mathbb{I}\rightarrow \mathbb{R}^3, \mathbb{I} \subset \mathbb{R}$. 
The derivative $\dot \vec{x}(t_0)=\frac{d\vec x}{dt}|_{t=t_0}$ is called the velocity vector. 
The derivative $\ddot \vec x(t_0)=\frac{ d^2 \vec x}{dt^2}|_{t=t_0}$ is called the acceleration vector. 
The image of motion (curve in $\mathbb{R}^3$) called trajectory and  curve in 
$\mathbb{T} \times \mathbb{R}^3$ which appears  as the graph of a motion, is called a world line.

The initial state of a mechanical system (the totality of positions and
velocities of its points at some moment of time) uniquely determines all of
its motion. In particular, the initial positions and velocities determine the accelerations.
This is called `Newton's principle of determinacy'.

\subsection{Newtonian mechanics}
\label{newt-mech}

Let us consider the mechanical system consisting of $n$ particles 
with masses $m_j,j=1\ldots n$ and word lines $\vec x_j(t)$.
Suppose that $n$ vector functions  
$\vec F_j:\mathbb{R}^{3 n}\times\mathbb{R}^{3 n}\times\mathbb{R}\rightarrow\mathbb{R}^3,j=1\ldots n$ 
are defined. Then the motion of all particles is determined by equations
\begin{equation}
  \label {newtoneq}
    m_j \ddot \vec x_j=\vec F_j(\vec x_k,\dot \vec x_k,t) \quad j,k=1\ldots n 
\end{equation}
and initial conditions
\begin{equation}
\label {newtonic}
   \vec x_j(t_0)=\vec y_j\quad\dot \vec x_j(t_0)=\vec v_j  \quad j=1\ldots n \end{equation}

Newton's equation \eqref{newtoneq} is  the basis of mechanics. Functions $\vec F_j$
are called forces they should be determined  for each specific mechanical system. If our system
is isolated i.e. includes explicitly all interacting particles, then galilean invariance puts some restrictions on forces.
Namely 
\begin{equation}
  \label {newtoneq-isolated}
    m_j \ddot \vec x_j=\vec F_j(\vec x_j-\vec x_k,\dot \vec x_j-\dot \vec x_k) \quad j,k=1\ldots n 
\end{equation}
practically such systems occur in celestial mechanics. In most other cases the system 
considered is not isolated and forces of general kind \eqref{newtoneq} appear as a result 
of interaction with environment.

Mathematics of  Cauchy problem (\ref{newtoneq},\ref{newtonic}) 
one can find in e.g. \cite{Arnold-ode}. Generally the unique solution exists if
all functions are smooth enough. However smoothness may be violated in physically
important cases, anyway functions $\vec x_j(t)$ should be continuous, $\vec v_j(t)$  may
have finite jumps as a results of strong force acting for very short time (e.g. 
elastic collision).

Note that initial conditions have nothing to do with causality. Causality makes
sense in Real world only, in Ideal world we are free to set final conditions or
fix $\vec x_j(t_1)$ and $\vec v_j(t_1)$ for some intermediate time $t_1$. In
many cases we are interested in solution of equation \eqref{newtoneq} with
fixed initial velocities and final coordinates  (aiming) or fixed coordinates
on both ends (variational problem). Anyway when conditions are set at two
different times  solution may be not unique (unless force is changing
relatively slow) \footnote {I discovered this fact around 1980 and was very
proud of myself until I found papers \cite{Maslov-1} and \cite{Maslov-2}.}.

\subsubsection{Constraints} 

In addition to Newton's equations the motion of the system may be restricted by
additional equations or inequalities. Constraints introduce two types of
difficulties in the solution of mechanical problems. First, the coordinates
$\vec x_j$ are no longer all independent, since they are connected by the
equations of constraint; hence the equations of motion  are not all
independent. Second, the forces of constraint  are not furnished a priori. They
are among the unknowns of the problem and must be obtained from the solution we
seek. Indeed, imposing constraints on the system  is simply another method of
stating that there are forces present in the problem that cannot be specified
directly but are known rather in terms of their effect on the motion of the
system. There is no general method of solving (even numerically) these
problems, they are tackled on case by case basis.

For the moment we should restrict ourselves with so called holonomic
constraints which have the form:
\begin{equation}
        f(\vec x_j,t)=0 \label {hol-const}
\end{equation}

\subsubsection{Numerical solutions}

Generally Newton's equations are to be solved numerically which is a difficult
task. Two large class of methods are in use: 
\begin{enumerate} 
    \item "Usual" one-time integrators when differential equations are replaced
        with finite difference ones defining snapshots at discrete time steps.
        These integrators allow accurate description of standard problem:
        initial conditions, finite time interval. 

\item Geometric integrators which are less accurate for any single trajectory
    but accurately present qualitative picture for large set of initial data
        and long times. Such integrators may calculate  coordinates and momenta
        at different times
\end{enumerate}
One can look for details in \cite{Hairer-Lubich-Wanner} for example. 

If we allow only potential forces i.e. 
\begin{equation}
  \label {potforce}
        \vec F_j(\vec x_k,t)=-\frac{\partial U(\vec x_k,t)}{\partial \vec x_j} \quad j,k=1\ldots n 
\end{equation}
we shall see that equations of motion (\ref{newtoneq}) can be written as 
\begin{equation}
   \label {lagr}
   \begin{split}
      \vec p_j&=\frac{\partial L(\vec x_k,\vec v_k,t)}{\partial \vec v_j} \\ 
        \dot \vec p_j&= \frac{\partial L(\vec x_k,\vec v_k,t)}{\partial \vec x_j}\quad j,k=1\ldots n 
   \end{split}
\end{equation}  
where 
\begin{equation}
  \label{lagr3}
        L(\vec x_k,\vec v_k,t)=\Sigma_1^n \frac{m_k}{2} \vec v_k \cdot \vec v_k-U(\vec x_k,t) 
\end{equation}

If we can solve the equation 
\[ \vec p_j=\frac{\partial L(\vec x_k,\vec v_k,t)}
{\partial \vec v_j} \] 
for $ \vec v_j $
which of course is not always possible, we can write the equations \eqref{lagr}
in more symmetrical way:
\begin{equation}
   \label {ham}
   \begin{split}
        \dot \vec x_j&=\frac{\partial H(\vec x_k,\vec p_k,t)}{\partial \vec p_j}  \\ 
        \dot \vec p_j&= -\frac{\partial H(\vec x_k,\vec p_k,t)}{\partial \vec x_j}\quad j,k=1\ldots n    
   \end{split}
\end{equation}
where 
\begin{equation}
        H(\vec x_k,\vec p_k,t)=\Sigma_1^n \frac{1}{2 m_k} \vec p_k \cdot \vec p_k+U(\vec x_k,t) \label{ham3}
\end{equation}
Note that functions $L$ (Lagrangian) and $H$ (Hamiltonian) are connected by Legendre transform:
\begin{equation}
        L(\vec x_k,\vec v_k,t)=\Sigma_1^n  \vec p_k \cdot \vec v_k -H(\vec x_k,\vec p_k,t)|_{\vec p_k=m_k\vec v_k}
\quad k=1\ldots n \label {lagr-ham}
\end{equation}


\subsection{Principle of least action}
\label{least-action}
\begin{definition}
        The equation
        \[\frac{d}{d t}(\frac{\partial L}{\partial \dot x})-\frac{\partial L}{\partial  x}=0 \]
        is called the Euler-Lagrange equation for the functional
        \[\Phi=\int_{t_1}^{t_2} L(x,\dot x,t) dt \]
\end{definition}
Let \( \gamma= \vec x(t) \quad t_1 \le t \le t_2 \) be a curve  in  $\mathbb{R} \times \mathbb{R}^n$ and 
$L: \mathbb{R}^n \times \mathbb{R}^n \times \mathbb{R}\rightarrow \mathbb{R}$ a function of $2 n +1$
variables, then
\begin{theorem}
         The curve $\gamma$ is an extremal of the functional \[\Phi=\int_{t_1}^{t_2} L(\vec x,\dot \vec x,t) dt \]
        on the space of curves joining $(t_1, \vec x_1)$ and $(t_2,\vec x_2)$, if and only if the Euler-
Lagrange equation is satisfied along $\gamma$.
\end{theorem}
This is a system of n second-order equations, and the solution depends on
$2 n$ arbitrary constants. The $2 n$ conditions $x(t_1) = x_0$ and $x(t_2) = x_2$ are used
for finding them.
Two important remarks:
\begin{enumerate}
        \item There may be many extremals connecting two
given points or none at all.
\item The condition for a curve  to be an extremal of a functional does not depend
on the choice of coordinate system. This paves way to geometrisation of mechanics.
\end{enumerate}

It is easy to see that equation of motion (\ref{lagr}) are actually  Euler-Lagrange equations 
for lagrangian \eqref{lagr3}. 
So we have  Hamilton's principle of least action.

\section{Deeper Insight}
\label{cm:di}

Classical mechanics is usually presented as the most settled part of
physics: three laws, a variational principle, and a semester of
worked examples. By the time a student reaches quantum theory or
general relativity, mechanics has already receded into the
background as a body of technique rather than a live subject. This
impression is not wrong, exactly, but it is incomplete. If one
presses on a few of the concepts that the standard course treats as
self-evident --- \emph{force}, \emph{constraint}, \emph{system} ---
the apparatus turns out to rest on assumptions that are easy to
state and easy to violate. What happens is instructive: relaxing a
single tacit assumption does not merely complicate the bookkeeping,
it can change the physics outright, produce equations that look
paradoxical on first encounter, or reveal that two "obviously
equivalent" formulations of mechanics are not equivalent at all.


I believe that least action principle  is nothing but (powerful) heuristic. Indeed, to compute 
action functional we should define explicitly {\em all} curves connecting given pair of points.
This seems to be impossible or at least very difficult problem. Fortunately we do not need it.
All what we really need is to set configuration space, choose Lagrange function 
\eqref{lagr3} and get equations of motion \eqref{lagr}. It does not matter at all
if we have minimum, maximum, saddle point or whatever property of some functional.

On the other hand having equations of motion as system of ordinary differential equations,
we can transform it into equivalent one, say by multiplying every equation by arbitrary 
function which never becomes zero. This new system may correspond to another Lagrangian and
the difference between new and old one need not come down to total derivative. Most probably
new system will not be Lagrangian at all.

Let us look at two-dimensional harmonic oscillator
\footnote{This is a simple example of difference between the
symmetry of equations and that of Lagrangian, see \ref{sp:eq} }.

Usual Lagrangian is
\begin{equation}
   \label {ho1}
    L_1= \frac{1}{2} (\dot x^2+\dot y^2) -\frac{ \omega^2}{2} ( x^2+ y^2)
\end{equation}
which yields equation  of motions:
\begin{equation}
   \label {hoe1}
   \begin{split}
       \ddot x+ \omega^2 x & =0 \\
       \ddot y+ \omega^2 y & =0 
   \end{split}
\end{equation}
Same equations of motion come form 
another Lagrangian
\begin{equation}
   \label {ho2}
    L_2= \frac{1}{2} (\dot x^2-\dot y^2) -\frac{ \omega^2}{2} ( x^2- y^2)
\end{equation}
or
\begin{equation}
   \label {ho3}
    L_3= \dot x\dot y - \omega^2  x y
\end{equation}
corresponding  Hamiltonians 
\begin{equation}
   \label {ho2h}
    H_2= \frac{1}{2} ( p_x^2-p_y^2) +\frac{ \omega^2}{2} ( x^2- y^2)
\end{equation}
and
\begin{equation}
   \label {ho3h}
    H_3= p_x p_y +\omega^2  x y
\end{equation}
are not even positive definite!

\section{Gauge invariance} 
\label{cm:gi}

From now on we shall use generalized coordinates $ q=\{q_j, j=1\ldots n\} $, where $ n $ is the number of
degrees of freedom. These coordinates may have or have not physical meaning.

We are used to think that gauge invariance is about electromagnetism, actually it appears
already in classical mechanics. It is easy to prove that adding the total derivative
\begin{equation}
   \label {tot-der}
    \frac{d \Phi (q,t)}{d t}=\frac{\partial \Phi(q,t)}{\partial t} + \frac{ \partial \Phi (q,t)}{\partial q_j} \dot q_j
\end{equation}
to Lagrangian does not change equations of motion. However canonical momenta
\begin{equation}
   \label {can-mom}
   p_j=\frac{ \partial L(q,\dot q,t)}{\partial \dot q_j}
\end{equation}
becomes \[ P_j=p_j+\frac{\partial \Phi}{\partial q_j} \] Transformation $ p,q \rightarrow P,q $ is evidently
canonical and new Hamiltonian becomes 
\[ H_{new}=H(p(P),q)- \frac{\partial \Phi}{\partial t}\]
The fact that canonical momenta have no physical meaning is not a weak but strong point of Hamiltonian 
mechanics.

I skip the  discussion of standard things like Poisson brackets, integrals of motion and all that.

\section{"Spin" in constant magnetic field}
\label{cm:mm}
We used to natural systems where hamiltonian is quadratic form over momenta,
but there are physical systems with hamiltonians linear over momenta.
So lagrangian does not exist. Let us consider 
dynamic of point magnetic moment in constant magnetic field has not only
clear physical meaning but also significant practical value. I investigated
this problem (actually slightly more general one) in connection with control of MRI
scanner \cite{RPZ4},\cite{RPZ3},\cite{RPZ2},\cite{RPZ1}.


Hamiltonian is $ H=\vec M \cdot \vec B $ where $\vec M$ is magnetic moment 
and $ \vec B $ is constant magnetic field.
It is convenient to write $ \vec M=\frac{e}{m c}\vec S $ where $ \vec S $
is angular moment ("spin")  and  include gyromagnetic ratio into magnetic field
$ \vec \Omega = \frac{e}{m c}\vec B$. Components of $ \vec S $ are not
canonical variables, their Poisson brackets are
\begin{equation}
   \label {pb}
    \{S_i,S_j\}= \epsilon_{ijk} S_k
\end{equation}
However they define phase space (sphere)  and equation of motion:
\begin{equation}
   \label {sp1}
    \dot S_i=\{S_i,H\}
\end{equation} 
or 
\begin{equation}
   \label {sp2}
    \dot {\vec S} =-\vec S \times \vec \Omega
\end{equation}

In canonical variables 
\begin{equation}
   \label {sp3}
   \begin{split}
       S_x & = \sqrt{S^2-p^2} \cos q \\
       S_y & = \sqrt{S^2-p^2} \sin q \\
       S_z & = p,\quad -S\le p \le S
   \end{split}
\end{equation}
Let $ \vec \Omega =(0,0, \Omega) $ then $ H=p \Omega $ and canonical equations are:
\begin{equation}
   \label {sp4}
   \begin{split}
       \dot p & =0\\
       \dot q & = \Omega
   \end{split}
\end{equation}
Lagrangian $ L=0 $, so we can not get equations of motion as second order equation.

Moment $ p $ is integral of motion which is fine but velocity is not connected to 
moment at all, it is constant. For pair of points in phase space generally there is
no trajectory passable for given time. Variation problem simply can not be posed.

 Hamiltonian $ H=p \Omega $ looks like that for harmonic oscillator in action-angle variables
and so we can  transform it to usual quadratic form, but the condition $  -S\le p \le S$
makes this transformation pure formal.

